\section{Cardinal-characteristic witness constructions}
\label{annex:p06-cardinal-characteristics}

\begin{annexitem}{annex:p06.cc.ideal-bounds}{thm:p06-cc-ideal-bounds}{Ideal witnesses and regular additivity}
\begin{proof}
A cover of \(X\) has union outside the proper ideal, giving
\(\operatorname{add}(I)\leq\operatorname{cov}(I)\). The singletons of any
\(A\notin I\) give \(\operatorname{add}(I)\leq\operatorname{non}(I)\).
A cofinal base covers \(X\), because each singleton is contained in one
base member. If \(\mathcal B\) is cofinal, choose \(x_B\in X\setminus B\)
for every \(B\in\mathcal B\). The set of chosen points is not in \(I\):
otherwise it would be included in some \(B\in\mathcal B\), contradicting
the choice of \(x_B\). This proves the remaining two bounds.

Countable unions stay in a \(\sigma\)-ideal, so its additivity is
uncountable. Put \(\kappa=\operatorname{add}(I)\). If
\(\lambda=\operatorname{cf}(\kappa)<\kappa\), partition a witnessing
\(\kappa\)-sized family into \(\lambda\) pieces of size less than
\(\kappa\). Each piece has union in \(I\) by minimality; the union of
those \(\lambda\) unions is again in \(I\), a contradiction. Thus
\(\kappa\) is regular. Finally the Borel subsets of a second-countable
real space are coded by countable labelled trees. A basic open has its
one-node code; a complement adds a new root, and a countable union adds
a root above countably many disjoint copied trees. Thus the coded sets
form a sigma algebra containing the countable base and hence all Borel
sets; the reverse inclusion follows by well-founded evaluation of each
code. The countable trees on finite sequences of naturals, with countable
node labels, have at most \(\mathfrak c\) possible codes. A Borel
cofinal base therefore has size at most \(\mathfrak c\).
\end{proof}
\end{annexitem}

\begin{annexitem}{annex:p06.cc.combinatorial}{thm:p06-cc-combinatorial}{Diagonal bounds and alternating intervals}
\begin{proof}
A countable family \(\{f_i:i<\omega\}\) is bounded by
\(g(n)=1+\max_{i\leq n}f_i(n)\). A dominating family cannot be bounded:
if \(g\) bounded it, domination of \(g+1\) would contradict that bound.
This gives \(\aleph_1\leq\mathfrak b\leq\mathfrak d\), and all of
\(\omega^\omega\) is a dominating family of size \(\mathfrak c\).
Regularity of \(\mathfrak b\) follows by partitioning an unbounded family
of size \(\mathfrak b\) into fewer than \(\mathfrak b\) smaller pieces:
bound each piece and then bound the smaller family of bounds. Transitivity
of \(\leq^*\) gives the contradiction.

No countable splitting family \(\{S_i:i<\omega\}\) suffices. Starting
with \(A_0=\omega\), choose recursively an infinite
\(A_{i+1}\subseteq A_i\) wholly in \(S_i\) or wholly in its complement.
Choose increasing \(a_i\in A_i\); the infinite set \(\{a_i:i<\omega\}\)
is almost contained in the chosen side for every \(S_i\), and is split
by none of them.

For \(f\in\omega^\omega\), set \(a_0=0\) and
\(a_{n+1}=1+\max\{a_n,f(a_n)\}\), and let
\(S_f=\bigcup_n[a_{2n},a_{2n+1})\). For infinite \(A\), let
\(g_A(k)=\min(A\setminus(k+1))\). If \(g_A\leq^*f\), then every
sufficiently late interval \([a_n,a_{n+1})\) meets \(A\), because
\(a_n<g_A(a_n)\leq f(a_n)<a_{n+1}\). Thus \(S_f\) splits \(A\).
A dominating family supplies a splitting family, proving
\(\mathfrak s\leq\mathfrak d\).

If \(\mathcal A\) is a family of infinite sets with
\(|\mathcal A|<\mathfrak b\), bound \(\{g_A:A\in\mathcal A\}\) by
one \(f\). The same \(S_f\) splits every member, so
\(\mathfrak b\leq\mathfrak r\). All infinite subsets together form
a reaping family: if \(S\) is infinite use \(A=S\), and if \(S\)
is finite use \(A=\omega\setminus S\). Hence
\(\mathfrak r\leq\mathfrak c\).
\end{proof}
\end{annexitem}

\begin{annexitem}{annex:p06.cc.real-ideals}{thm:p06-cc-real-ideals}{Borel bases and incomparable smallness}
\begin{proof}
For the Baire input used here and below, recall the complete-metric-space
argument. Given dense open \(G_n\) in a nonempty complete metric space
and a nonempty open \(U\), choose nested nonempty closed balls, the
first inside \(U\cap G_0\), the next inside the interior of the previous
ball and \(G_1\), and so on, with radii tending to zero. The centers
are Cauchy and their limit lies in every ball by closedness and completeness.
It belongs to \(U\cap\bigcap_nG_n\). If an isolated point occurs,
the singleton ball already lies in every dense open set and completes
the argument. Thus a countable closed cover of a nonempty complete
space has a member with nonempty relative interior, the form used below.
Countable subadditivity gives a \(\sigma\)-ideal of outer-measure-zero
sets; singletons are null and \([0,1]\) is not null. Every null set is
contained in a null \(G_\delta\): choose open supersets of measure
less than \(2^{-n}\) and intersect them. For category, the closure of
a nowhere dense set is closed nowhere dense. Thus a meager set is
contained in a meager \(F_\sigma\). The union of countably many such
countable covers is again meager; singletons are nowhere dense.
Baire's theorem says that \(\mathbb R\) is not meager. These facts
give the required proper ideals and Borel bases, and the preceding ideal
bounds give all four cardinal bounds.

Enumerate \(\mathbb Q=\{q_j:j<\omega\}\). For every \(n\), choose
positive radii \(\epsilon_{nj}\) with
\(\sum_j2\epsilon_{nj}<2^{-n}\) and put
\(U_n=\bigcup_j(q_j-\epsilon_{nj},q_j+\epsilon_{nj})\).
Each \(U_n\) is open dense and has measure at most \(2^{-n}\).
Their intersection \(E\) is a dense \(G_\delta\), null and comeager.
It is not meager, since otherwise \(E\) and its meager complement
would make \(\mathbb R\) meager. The complement is not null, since
it is conull. This gives witnesses to both failed ideal inclusions,
not merely a comparison of the two cardinal diagrams.
\end{proof}
\end{annexitem}

\begin{annexitem}{annex:p06.cc.category-domination}{thm:p06-cc-category-domination}{Eventually bounded cones}
\begin{proof}
In \(\omega^\omega\), the set
\(C_{g,N}=\{f:(\forall n\geq N)f(n)\leq g(n)\}\) is closed nowhere
dense: extend any finite stem at an unused coordinate \(n\geq N\)
with value \(g(n)+1\). Consequently
\(C_g=\{f:f\leq^*g\}=\bigcup_N C_{g,N}\) is meager. Any family
of size less than \(\mathfrak b\) is contained in such a cone; a
dominating family gives a cover by such cones.

For completeness on the stated real-line domain, the following direct
code avoids importing a homeomorphism lemma. Enumerate the rationals
without repetitions as \(\{q_n:n<\omega\}\). For irrational \(x\),
put \(f_x(n)=\min\{k\geq1:1/k\leq|x-q_n|\}\); the minimum exists
because the distance is positive. For positive-valued \(g\), set
\[
 K_{g,N}=\{x\in\mathbb R:(\forall n\geq N)|x-q_n|\geq1/g(n)\}.
\]
This set is closed nowhere dense. Indeed, the tail rationals are dense,
and any open interval contains a \(q_n\), \(n\geq N\), and a smaller
open neighborhood of that point disjoint from \(K_{g,N}\).
If \(f_x\leq^*g\), then eventually
\(|x-q_n|\geq1/f_x(n)\geq1/g(n)\); hence
\(x\in\bigcup_NK_{g,N}\).

For a real \(X\) of size below \(\mathfrak b\), bound all
\(f_x\), \(x\in X\setminus\mathbb Q\), by one \(g\), replacing
\(g(n)\) by \(\max\{1,g(n)\}\). Then \(X\) is contained in
the meager set \(\mathbb Q\cup\bigcup_NK_{g,N}\), proving the
uniformity bound. Normalize every member of a dominating \(D\) in
the same way. The meager sets
\(\mathbb Q\cup\bigcup_NK_{g,N}\), \(g\in D\), cover all reals,
because every irrational has its \(f_x\) eventually bounded by a
member of \(D\). Their number is at most \(\mathfrak d\), proving
the covering bound.
\end{proof}
\end{annexitem}

\begin{annexitem}{annex:p06.cc.category-coding}{thm:p06-cc-category-coding}{Closed nowhere dense witnesses recover eventual bounds}
\begin{proof}
This is the closed-witness construction of Fremlin 522H, with the choices
and the eventual, rather than pointwise, diagonal bound made explicit.
Enumerate a base of nonempty bounded rational intervals as
\(\{U_k:k<\omega\}\). Inside each \(U_k\), choose pairwise disjoint
nonempty open intervals \(V_{kl}\), and enumerate the rationals in each
as distinct \(\{x_{kli}:i<\omega\}\). Put
\[
 E_k(f)=\{x_{kli}:l<\omega,\ i\leq f(l)\}.
\]
This is nowhere dense. If an open interval meets no \(V_{kl}\), it
misses \(E_k(f)\); otherwise restrict it inside one \(V_{kl}\), where
only finitely many selected points occur, and then avoid those points.

Construct a sequence of protective intervals \(W_n\) and a possibly
empty selection of sets \(E_k(f)\) at each stage. At stage \(n\),
choose a nonempty basic interval \(W_n\) whose closure is contained
in \(U_n\) and avoids the closures of the finitely many previously
selected \(E_k(f)\). If \(U_n\) has points outside
\(\bigcup_{i<n}\overline{W_i}\), choose \(W_n\) sufficiently small
that \(U_n\setminus\bigcup_{i\leq n}\overline{W_i}\) still has a
nonempty open subset. These requirements are compatible: the old
selected sets are nowhere dense, and one can leave a second small open
interval outside the finite closed union whenever that complement is
nonempty. If the new complement is nonempty, choose a basic
\(U_{k_n}\) inside it and select \(E_{k_n}(f)\); otherwise select no
set at this stage. At stage zero the complement is nonempty, so at least
one set is selected. Define \(\Phi(f)\) to be the closure of the union
of all selected sets.

Each selected set is disjoint from every \(\overline{W_n}\): this
holds for past selections by the choice of \(W_n\), and for present
and future selections by the protective finite closed unions. Hence
the open \(W_n\) misses \(\Phi(f)\), and lies in \(U_n\).
Every basic interval therefore contains an open interval outside
\(\Phi(f)\); the latter is closed nowhere dense and nonempty.
Moreover, if an open \(G\) meets \(\Phi(f)\), take
\(U_n\subseteq G\) meeting it. Then \(U_n\) meets some selected
set, whose point lies outside \(\bigcup_{i\leq n}\overline{W_i}\).
The complement at stage \(n\) is thus nonempty, so a set was selected
at that stage, and \(E_{k_n}(f)\subseteq G\cap\Phi(f)\).

For meager \(A\), choose closed nowhere dense \(F_j\) covering it.
For every \(j,k\), define
\(g_{jk}(l)=\min\{i:x_{kli}\notin F_j\}\). The minimum exists
because an open interval is not contained in \(F_j\), and its rational
points are dense. Choose \(\Psi(A)\) eventually dominating all the
countably many \(g_{jk}\), by the diagonal construction already proved.
If \(\Phi(f)\subseteq A\), Baire's theorem on the nonempty closed
space \(\Phi(f)\) gives \(j\) and an open \(G\) with
\(\varnothing\ne G\cap\Phi(f)\subseteq F_j\). The preceding local
property gives an \(E_k(f)\) inside this set. Since all
\(x_{kli}\) with \(i\leq f(l)\) belong to \(F_j\),
\(f(l)<g_{jk}(l)\) for every \(l\). Therefore
\(f\leq^*\Psi(A)\), as required.

For an unbounded \(F\) of size \(\mathfrak b\), the union of
\(\{\Phi(f):f\in F\}\) cannot be meager, since a meager container
would bound every \(f\) by the same \(\Psi(A)\). This proves the
additivity inequality. If \(\mathcal B\) is a cofinal meager base,
then for every \(f\) some \(B\in\mathcal B\) contains \(\Phi(f)\);
thus \(\{\Psi(B):B\in\mathcal B\}\) is dominating, proving the
cofinality inequality.
\end{proof}
\end{annexitem}

\begin{annexitem}{annex:p06.cc.category-identities}{thm:p06-cc-category-identities}{One translation and one eventual bound}
\begin{proof}
Enumerate the rationals as \(\{q_n:n<\omega\}\) with every rational
repeated infinitely often. For \(f\in\omega^\omega\), put
\[
 L_f=\bigcap_N\bigcup_{n\geq N}
       (q_n-2^{-f(n)},q_n+2^{-f(n)}),\qquad E_f=\mathbb R\setminus L_f.
\]
Every tail union is open dense, so \(E_f\) is meager and disjoint from
\(\mathbb Q\). If \(f\leq^*g\), tail interval inclusion gives
\(E_f\subseteq E_g\). If closed sets \(H_j\) all miss \(\mathbb Q\),
choose \(f(n)\) so the interval around \(q_n\) misses
\(\bigcup_{j\leq n}H_j\). Then \(\bigcup_jH_j\subseteq E_f\):
a point in \(H_j\) lies in none of the intervals from index \(j\) on.

Let \(\mathcal A\) be a family of meager sets of size less than both
\(\mathfrak b\) and \(\operatorname{cov}(\mathcal M)\).
For each \(A\) choose a closed nowhere dense cover with union \(F_A\).
The sets \(F_A+\mathbb Q\) are meager; their union cannot cover
\(\mathbb R\). Choose \(x\) outside it. Every closed piece of
\(F_A-x\) misses \(\mathbb Q\), so choose \(f_A\) with
\(A-x\subseteq E_{f_A}\). Bound all \(f_A\) eventually by one \(g\).
Then \(\bigcup\mathcal A\subseteq x+E_g\), a meager set. This proves
the lower bound for additivity. The opposite bound follows from the ideal
square and category-coding theorem.

For cofinality, choose a nonmeager \(X\) of size
\(\operatorname{non}(\mathcal M)\) and a dominating
\(D\) of size \(\mathfrak d\). For any meager \(A\), choose
\(x\in X\setminus(F_A+\mathbb Q)\). As before,
\(A\subseteq x+E_f\) for some \(f\); choose \(g\in D\) with
\(f\leq^*g\), obtaining \(A\subseteq x+E_g\). Therefore
\(\{x+E_g:x\in X,g\in D\}\) is cofinal, of size at most
\(\max\{|X|,|D|\}\). The reverse bound is the ideal square together
with the category-coding theorem.
\end{proof}
\end{annexitem}

\begin{annexitem}{annex:p06.cc.cross-covers}{thm:p06-cc-cross-covers}{Covering by translates of opposite smallness}
\begin{proof}
Use the null comeager \(E\) and meager conull \(H=\mathbb R\setminus E\)
constructed in the real-ideals proof. Let \(A\) be a nonmeager set of
least possible size. For every \(y\), \(y-E\) is comeager and hence
meets \(A\); otherwise \(A\) would be contained in its meager
complement. Thus \(y=a+e\) for some \(a\in A,e\in E\), and
the null sets \(\{a+E:a\in A\}\) cover \(\mathbb R\).
This proves the first inequality.

For the second take a nonnull \(A\) of least possible size. The set
\(y-H\) is conull and meets \(A\), because its complement is null
and cannot contain \(A\). Therefore the meager translates
\(\{a+H:a\in A\}\) cover \(\mathbb R\). Translation and reflection
preserve outer-measure-zero sets and meagerness; no measurability or
Baire property of \(A\) was used.
\end{proof}
\end{annexitem}

\begin{annexitem}{annex:p06.cc.null-localization}{thm:p06-cc-null-localization}{Independent events and summable interval codes}
\begin{proof}
We construct the two implications separately, following Fremlin 522M.
First choose open sets \(G_{ij}\subseteq(0,1)\), for \(i,j<\omega\),
of Lebesgue measure \(2^{-i}\), independent in every finite subfamily.
Here is an explicit measure construction. Enumerate the pairs with
\(i\geq1\), assign pairwise disjoint finite blocks of \(i\) binary
digit positions to them, and require all digits in the assigned block
to be zero. Each event is a finite union of dyadic half-open intervals;
take its interior. Only finitely many endpoints are removed from each
event, so its measure and all finite intersection measures are unchanged.
Disjoint digit blocks give the required product probabilities.
More explicitly, for a finite collection take \(L\) past every specified
digit position. Exactly \(2^{L-\sum i}\) of the \(2^L\) binary intervals
of length \(2^{-L}\) satisfy the specified zeros; their total measure
is \(2^{-\sum i}\), the product of the individual probabilities.
For \(i=0\), use \(G_{0j}=(0,1)\); these probability-one events do
not affect independence. Define
\[
 a(f)=\bigcap_N\bigcup_{i\geq N}G_{i,f(i)}.
\]
The measure of a tail union is at most \(\sum_{i\geq N}2^{-i}\),
which tends to zero. Thus \(a(f)\) is null.

For null \(A\subseteq\mathbb R\), choose a positive-measure compact
\(K\subseteq(0,1)\setminus A\) such that every nonempty relatively
open part of \(K\) has positive measure. To make this choice, cover
\(A\) by an open \(U\) with \(\mu(U)<1/4\). Then
\(K_0=[1/4,3/4]\setminus U\) is compact, avoids \(A\), and has
measure at least \(1/2-\mu(U)>0\). Remove from \(K_0\) the union
of all intersections with basic open intervals \(V\) for which
\(\mu(K_0\cap V)=0\). This is a relatively open null set, so the
remaining \(K\) is compact of the same positive measure. If a basic
interval meets \(K\), its intersection with \(K_0\) has positive
measure, and removing the null union leaves positive measure in \(K\).
This gives the stated support property without an inner-regularity lemma.
Let \(W_n\) be a
countable base of its nonempty relatively open parts. Set
\[
 I_{ni}=\{j:W_n\cap G_{ij}=\varnothing\}.
\]
For fixed \(n\),
\(\sum_i2^{-i}|I_{ni}|<\infty\). Indeed, \(W_n\) is disjoint
from every event indexed by these pairs. If their sum of probabilities
diverged, finite independence would give
\[
 \mu\Bigl(\bigcap_{(i,j)\in F}((0,1)\setminus G_{ij})\Bigr)
 =\prod_{(i,j)\in F}(1-2^{-i})
 \leq\exp\Bigl(-\sum_{(i,j)\in F}2^{-i}\Bigr).
\]
Increasing finite \(F\) would force \(\mu(W_n)=0\), a contradiction.
This also proves each \(I_{ni}\) is finite.
Choose \(k_n\) so \(2^{-i}|I_{ni}|\leq2^{-n-1}\) for \(i\geq k_n\),
and put \(b(A)(i)=\bigcup_{n:k_n\leq i}I_{ni}\).
Its width is at most \(2^i\sum_n2^{-n-1}=2^i\); hence it is a slalom.

If \(a(f)\subseteq A\), then \(K\) misses \(a(f)\) and is covered
by the relatively closed sets
\(K\setminus\bigcup_{i\geq m}G_{i,f(i)}\), \(m<\omega\).
Baire's theorem on the nonempty compact \(K\) gives a basic \(W_n\)
disjoint from one such tail union. For every \(i\geq m\), therefore,
\(f(i)\in I_{ni}\), and for \(i\geq\max\{m,k_n\}\) it belongs
to \(b(A)(i)\). This is the first required implication.

For the reverse direction, enumerate for every \(n\) the countable
family of finite unions of bounded rational open intervals of total
measure at most \(4^{-n}\), as \(\{H_{nj}:j<\omega\}\).
Every null \(A\) has a code \(c(A)\in\omega^\omega\) such that
\[
 A\subseteq\bigcap_N\bigcup_{n\geq N}H_{n,c(A)(n)}.
\]
To prove existence, for each \(m\) cover \(A\) by rational open
intervals whose summed lengths are at most \(2^{-m-1}\).
For \(A=\varnothing\), simply take the empty interval union at every
coordinate; otherwise the row covers supply infinitely many occurrences.
Enumerate all these interval occurrences in one sequence \(J_i\).
The total length is at most one, and each point of \(A\) is in
infinitely many occurrences, one from each row. Choose increasing
\(k(n)\) with \(k(0)=0\) such that the tail length from \(k(n)\)
is at most \(4^{-n}\). The finite union
\(\bigcup_{k(n)\leq i<k(n+1)}J_i\) has the required small measure
and is one of the \(H_{nj}\). Every point belongs to infinitely many
of these blocks, giving the asserted code.

For a slalom \(S\), define
\[
 e(S)=\bigcap_N\bigcup_{n\geq N}\bigcup_{j\in S(n)}H_{nj}.
\]
The measure of its tail cover is at most
\(\sum_{n\geq N}2^n4^{-n}\), tending to zero, so \(e(S)\) is null.
If \(c(A)(n)\in S(n)\) eventually, discarding finitely many code
coordinates does not change the limsup cover; thus \(A\subseteq e(S)\).
This proves the second implication and completes all four map constructions.
\end{proof}
\end{annexitem}

\begin{annexitem}{annex:p06.cc.meager-null-transport}{thm:p06-cc-meager-null-transport}{Finite-intersection open families and the containment bridge}
\begin{proof}
We first supply the open-family helper of Fremlin 522N, so it is not
hidden in a named localization lemma. In any nonempty space with a
countable \(\pi\)-base, for each finite \(t\) there is a countable
family \(\mathcal V_t\) of open sets such that every dense open set
contains a member and every intersection of at most \(t\) members
is nonempty. Start with \(\mathcal V_0=\{\varnothing\}\), using
the convention that an empty intersection is the whole nonempty space.
Let \(\mathcal U\) be the countable \(\pi\)-base closed under finite
unions. Given \(\mathcal V_t=\{H_i:i<\omega\}\), put
\[
 \mathcal J_i=\{J\subseteq i:\bigcap_{j\in J}H_j\ne\varnothing\},
\quad
 \mathcal V_{t+1}=\{U\cup H_i:U\in\mathcal U,
       \ (\forall J\in\mathcal J_i)\ U\cap\bigcap_{j\in J}H_j\ne\varnothing\}.
\]
For a dense open \(G\), choose \(H_i\subseteq G\), then choose
finitely many basic opens inside \(G\) meeting each nonempty open
intersection indexed by \(\mathcal J_i\); their union is an admissible
\(U\subseteq G\). To check the intersection property, take at most
\(t+1\) members \(U_k\cup H_{i_k}\), and select a largest index
\(i_*\). If it is repeated, the intersection of the distinct
\(H_{i_k}\) is nonempty by the induction hypothesis. If it is unique,
the intersection of the other at most \(t\) sets \(H_{i_k}\) is
nonempty and indexed by a subset of \(i_*\); the admissible
\(U_*\) meets it. Either way the whole displayed intersection is
nonempty. This completes the finite induction.

Enumerate a real \(\pi\)-base \(\{U_n:n<\omega\}\) with every
member occurring cofinally. Apply the helper inside \(U_n\) with
\(t=2^n\), and enumerate its countable family as
\(\{V_{ni}:i<\omega\}\). For meager \(A\), choose increasing
closed nowhere dense \(F_n\) covering it. Since \(U_n\setminus F_n\)
is dense open in \(U_n\), choose \(p(A)(n)\) with
\(V_{n,p(A)(n)}\cap F_n=\varnothing\).
For a slalom \(S\), define
\[
 O_n(S)=\bigcap_{i\in S(n)}V_{ni},\qquad
 q(S)=\mathbb R\setminus\bigcap_N\bigcup_{n\geq N}O_n(S).
\]
When \(S(n)=\varnothing\), set \(O_n(S)=U_n\).
Every \(O_n(S)\) is nonempty open in \(U_n\). Every tail union
is dense open, because every basic open recurs at arbitrarily late
indices. Thus \(q(S)\) is meager. If \(p(A)\in^*S\), take \(N\)
after which localization holds. For every fixed \(j\), all
\(O_n(S)\) with \(n\geq\max\{N,j\}\) avoid \(F_j\), since
they avoid \(F_n\supseteq F_j\). Hence \(F_j\subseteq q(S)\),
and \(A\subseteq q(S)\).

Now use the maps \(a,b\) from the preceding null-localization lemma:
set \(u(A)=a(p(A))\) and \(v(B)=q(b(B))\). If \(u(A)\subseteq B\),
then \(p(A)\in^*b(B)\), and consequently \(A\subseteq v(B)\).
This proves the containment bridge with all domains explicit.

For a family \(\mathcal A\) witnessing meager additivity, its mapped
null sets cannot have a null union \(B\): that would put every
\(A\in\mathcal A\) in the same meager \(v(B)\), contrary to the
witnessing union. Thus \(\operatorname{add}(\mathcal N)
\leq|\mathcal A|\). If \(\mathcal B\) is a cofinal null base, then
\(\{v(B):B\in\mathcal B\}\) is a cofinal meager base: for any
meager \(A\), choose \(B\) containing \(u(A)\), and the bridge
puts \(A\) in \(v(B)\). This gives the cofinality inequality.
\end{proof}
\end{annexitem}

\begin{annexitem}{annex:p06.cc.cichon}{thm:p06-cc-cichon}{Assembling the generating edges}
\begin{proof}
The endpoint bounds and the four same-ideal comparisons are the ideal
bounds and real-ideals theorems. The central \(\mathfrak b\leq\mathfrak d\)
edge is the combinatorial theorem. The category-coding theorem supplies
the two outer category/bounding edges, and the category-domination
theorem the other two. Translation supplies both crossing covering/
uniformity edges. Finally meager--null transport supplies additivity
and cofinality between the ideals. These are exactly the generators
displayed; transitivity of cardinal comparison supplies their consequences.
Each generator therefore has the earlier reciprocal construction, rather
than being assumed from the shape of a diagram.
\end{proof}
\end{annexitem}

\begin{annexitem}{annex:p06.cc.ch-collapse}{thm:p06-cc-ch-collapse}{The endpoint calculation under CH}
\begin{proof}
Every characteristic under discussion is at least \(\aleph_1\) and
at most \(\mathfrak c\), by the real-ideal and combinatorial bounds.
Under \(\mathsf{CH}\), these two endpoints are equal. The inequalities
therefore force each value to be \(\aleph_1\). This calculation supplies
no strict separation in a model with larger continuum.
\end{proof}
\end{annexitem}
